\documentclass[11pt,a4paper]{article}

\usepackage[utf8]{inputenc}
\usepackage[T1]{fontenc}
\usepackage{amsmath,amssymb,amsthm}
\usepackage{graphicx}
\usepackage{booktabs}
\usepackage{hyperref}
\usepackage{xcolor}
\usepackage[margin=25mm]{geometry}

\newtheorem{theorem}{Theorem}
\newtheorem{proposition}{Proposition}
\newtheorem{observation}{Observation}
\newtheorem{conjecture}{Conjecture}
\newtheorem{corollary}{Corollary}

\title{Perturbative and Non-Perturbative Wings of the Zeta Building:\\
Kreimer--Connes at $\zeta(s>0)$ and Vacuum Cosmology at $\zeta(s<0)$,\\
Unified by the Functional Equation}

\author{Wright Brothers Research Program}
\date{April 2026}

\begin{document}

\maketitle

\begin{abstract}
The Kreimer--Connes Hopf algebra framework established that
perturbative QFT renormalization produces motivic periods ---
specifically, multiple zeta values $\zeta(3), \zeta(5), \zeta(7),
\ldots$ --- as genuine physical content of loop integrals, not mere
regularization artifacts. We observe that this ``reification of
$\zeta$ at positive integers'' has a natural counterpart: the
vacuum zero-point energy, regularized via $\zeta$ at
\emph{negative} integers ($\zeta(-1), \zeta(-3), \zeta(-5), \ldots$,
equivalently Bernoulli numbers), also carries physical content.
The Riemann functional equation $\zeta(s) \leftrightarrow \zeta(1-s)$
connects these two sides. We propose that the \emph{full}
physical content of the $\zeta$ function comprises both wings:
the \textbf{perturbative wing} (Kreimer--Connes, $\zeta$ at
positive integers, loop integrals) and the
\textbf{non-perturbative wing} (vacuum cosmology, $\zeta$ at
negative integers, boundary-condition-dependent zero-point sums).
As evidence, we show: (i) the fine structure constant $\alpha^{-1}$
decomposes into an algebraic leading part from $\zeta(-1), \zeta(-3)$
(Bernoulli) plus transcendental corrections from $\zeta(6), \zeta(7)$
(periods) --- precisely the structure predicted by the
Deligne--Beilinson period conjectures; (ii) the dark energy density
$\Omega_\Lambda = 16|\zeta(-1)|\zeta(2)/\pi = 2\pi/9$ combines
both wings in a single observable, derivable independently from
Wheeler--DeWitt quantum cosmology with temporal Dirichlet--Neumann
boundaries; (iii) the Euler product truncation at $p=2$
($\zeta \to \zeta_{\neg 2}$), corresponding to DN boundary
conditions, is a non-perturbative selection rule invisible to
perturbation theory. This framework positions vacuum cosmology
as the natural complement --- not competitor --- to Kreimer--Connes
perturbative renormalization, connected by the functional equation.
\end{abstract}

\section{Introduction}

\subsection{Kreimer--Connes: $\zeta$ values are real in perturbative QFT}

The work of Kreimer~\cite{Kreimer1998} and
Connes--Kreimer~\cite{CK1998,CK2000,CK2001} established that
renormalization in perturbative quantum field theory is not an
ad hoc subtraction procedure but a precise algebraic operation:
the Birkhoff decomposition of characters on a Hopf algebra of
Feynman graphs. A key consequence is that the \emph{finite}
(renormalized) parts of Feynman integrals evaluate to
\textbf{motivic periods} --- algebraic combinations of multiple
zeta values (MZVs) and
polylogarithms~\cite{Brown2015,BrownKreimer2013}:
\begin{equation}
\text{Renormalized Feynman integral} \in
\mathbb{Q}[\zeta(3), \zeta(5), \zeta(7), \ldots, \mathrm{Li}_n(\tfrac{1}{2}), \ln 2, \pi^2].
\end{equation}

For example, the two-loop QED contribution to the electron
anomalous magnetic moment contains
$\frac{3}{4}\zeta(3)$~\cite{Laporta2017}. The three-loop involves
$\zeta(5)$. Higher loops bring $\zeta(7)$ and MZVs of increasing
depth.

Connes and Marcolli~\cite{ConnesMarcolli2008} further showed that
a \emph{cosmic Galois group} acts on these periods, giving the
arithmetic content of perturbative QFT a Galois-theoretic structure
analogous to the motivic Galois group of algebraic geometry.

The philosophical content of this program is profound:
\textbf{$\zeta$ values at positive integers are not regularization
artifacts but genuine physical invariants of quantum field theory.}
They are periods of algebraic varieties (complements of graph
hypersurfaces), and their appearance in physical observables
(anomalous dimensions, $\beta$-functions, $g-2$) is
mathematically inevitable.

\subsection{The missing wing}

The Kreimer--Connes framework operates within perturbation theory.
It addresses the arithmetic of \emph{loop integrals} --- fluctuations
\emph{around} a background vacuum. It does not address the vacuum
itself.

However, the vacuum also has arithmetic content. The zero-point
energy of a quantum field in a cavity of size $a$ is formally
\begin{equation}\label{eq:zpe}
E_0 = \frac{\hbar}{2}\sum_n \omega_n
\end{equation}
which diverges but, via $\zeta$-regularization, yields a finite
physical quantity --- the Casimir energy. The regularization involves
$\zeta$ at \emph{negative integers}:
\begin{equation}
\zeta(-1) = -\frac{1}{12} = -\frac{B_2}{2}, \quad
\zeta(-3) = \frac{1}{120} = -\frac{B_4}{4}, \quad \ldots
\end{equation}
These are \emph{Bernoulli numbers} --- rational, algebraic, on the
opposite side of the $\zeta$ function from the periods $\zeta(3),
\zeta(5), \ldots$ that Kreimer--Connes handles.

The question of this paper is:
\begin{quote}
\emph{If Kreimer--Connes establishes that $\zeta$ at positive
integers is physically real (in perturbative QFT), is $\zeta$ at
negative integers also physically real (in non-perturbative vacuum
physics)?}
\end{quote}

We argue \textbf{yes}, and show that the Riemann functional equation
provides the bridge.

\subsection{The functional equation as a physical duality}

The Riemann functional equation
\begin{equation}\label{eq:functional}
\zeta(s) = 2^s \pi^{s-1}\sin\!\bigl(\tfrac{\pi s}{2}\bigr)
\,\Gamma(1-s)\,\zeta(1-s)
\end{equation}
relates $\zeta(s)$ to $\zeta(1-s)$. At $s = 2$ and $s = -1$:
\begin{equation}
\zeta(2) = \frac{\pi^2}{6}\quad \longleftrightarrow \quad
\zeta(-1) = -\frac{1}{12}.
\end{equation}
The left side is a transcendental period; the right side is a
rational Bernoulli value. They are ``dual'' under the functional
equation.

We propose that this mathematical duality has a \emph{physical}
counterpart:

\begin{center}
\begin{tabular}{lll}
\toprule
 & \textbf{Perturbative wing} & \textbf{Non-perturbative wing} \\
\midrule
$\zeta$ arguments & Positive integers & Negative integers \\
Values & Periods ($\pi^{2k}$, MZVs) & Bernoulli ($B_{2k}/(2k)$) \\
Physics & Loop integrals & Vacuum zero-point \\
Framework & Kreimer--Connes & Wheeler--DeWitt DN \\
Observables & $a_e$, $\alpha$ running & $\Omega_\Lambda$, Casimir \\
Math nature & Transcendental & Algebraic (rational) \\
\bottomrule
\end{tabular}
\end{center}

The functional equation (\ref{eq:functional}) connects these wings,
and we will show that specific physical observables --- notably
$\alpha^{-1}$ and $\Omega_\Lambda$ --- involve \emph{both} wings
simultaneously.

\section{The two wings in detail}

\subsection{Perturbative wing: Kreimer--Connes}

We briefly review the Kreimer--Connes results relevant to our
argument.

\textbf{Hopf algebra of Feynman graphs.}
Let $\mathcal{H}$ be the commutative Hopf algebra generated by
one-particle irreducible Feynman graphs, with coproduct encoding
subdivergence structure~\cite{Kreimer1998}. A \emph{character}
$\phi: \mathcal{H} \to \mathcal{A}$ (where $\mathcal{A}$ is an
algebra of Laurent series in the regularization parameter
$\varepsilon$) encodes the evaluation of each graph.

\textbf{Birkhoff decomposition.}
The character $\phi$ admits a unique factorization
$\phi = \phi_-^{-1} \star \phi_+$, where $\phi_-$ contains the
poles (counterterms) and $\phi_+$ is holomorphic (renormalized
values)~\cite{CK2000}. This is the Bogoliubov--Parasiuk--Hepp--Zimmermann
(BPHZ) renormalization in algebraic form.

\textbf{Periods.}
The renormalized values $\phi_+(\Gamma)$ for primitive graphs
$\Gamma$ evaluate to periods of mixed Tate motives over
$\mathbb{Z}$~\cite{Brown2015}. Concretely, they are
$\mathbb{Q}$-linear combinations of products of $\zeta(n)$
($n \geq 2$ integer) and $\ln 2$. For massless propagators, only
odd $\zeta$-values ($\zeta(3), \zeta(5), \ldots$) appear at the
single-zeta level.

\textbf{Physical consequence.}
The QED electron $g-2$, computed to five loops, contains
$\zeta(3), \zeta(5)$ at 2--3 loops, and deeper MZVs at higher
loops~\cite{Laporta2017,Aoyama2020}. These are \emph{measured}
through the electron $g-2$ experiment, making them empirically
verified physical quantities.

\subsection{Non-perturbative wing: vacuum zero-point}

The vacuum zero-point energy (\ref{eq:zpe}) depends on
\emph{boundary conditions} imposed on the field modes:

\begin{itemize}
\item \textbf{Dirichlet--Dirichlet (DD)} or
\textbf{Neumann--Neumann (NN)}:
All-integer modes $n = 1, 2, 3, \ldots$. Regularized sum
$\sum n^{-s}|_{s=-1} = \zeta(-1) = -1/12$.

\item \textbf{Dirichlet--Neumann (DN)}:
Odd-integer modes $n = 1, 3, 5, \ldots$. Regularized sum
$\sum_{\text{odd}} n^{-s}|_{s=-1} = \zeta_{\neg 2}(-1) = +1/12$.
\end{itemize}

where
\begin{equation}
\zeta_{\neg 2}(s) \equiv (1 - 2^{-s})\zeta(s)
\end{equation}
is the $\zeta$ function with the $p = 2$ Euler factor removed.

\textbf{Key difference from Kreimer--Connes:} this is
\emph{not} a loop integral. It is the sum over all modes of the
\emph{background vacuum}, without perturbative expansion. The
$\zeta$ values at negative integers ($\zeta(-1), \zeta(-3), \ldots$)
emerge not from Feynman diagrams but from the
\emph{mode spectrum} of the vacuum itself.

\textbf{Physical consequence.} The 3D Casimir effect gives a
measurable force $\propto \zeta(-3) = 1/120$~\cite{Casimir1948},
confirmed experimentally~\cite{Lamoreaux1997}. Boyer's DN Casimir
prediction involves $\zeta_{\neg 2}(-3) = -7/120$~\cite{Boyer1974},
not yet experimentally verified.

\subsection{Why Kreimer--Connes cannot see the non-perturbative wing}

Perturbation theory expands around a \emph{fixed} vacuum. The
vacuum itself is the zeroth-order background. Its zero-point
energy is subtracted (normal ordering) or absorbed into the
cosmological constant. Kreimer--Connes, operating within
perturbation theory, systematically handles all loop corrections
\emph{above} the vacuum, but the vacuum energy itself lies
\emph{outside} the Hopf algebra of Feynman graphs.

This is precisely why perturbative QFT cannot solve the
cosmological constant problem: $\Lambda$ is a non-perturbative
vacuum property, not a loop correction.

\section{Evidence: $\alpha^{-1}$ spans both wings}

\subsection{Decomposition}

We have previously reported~\cite{wb-alpha} the arithmetic formula
\begin{equation}\label{eq:alpha}
\alpha^{-1} \approx \underbrace{\frac{2}{|B_2|} + \frac{4}{|B_4|}}_{\text{algebraic (Bernoulli)}}
+ \underbrace{g \cdot 132 + 4(\zeta(6) - \zeta(7))}_{\text{transcendental (periods)}}
\end{equation}
where $2/|B_2| = 12$, $4/|B_4| = 120$, $g \approx 0.038$,
and $4(\zeta(6) - \zeta(7)) \approx 0.036$.

We now observe that this decomposition has a natural interpretation
in terms of the two wings:

\begin{itemize}
\item \textbf{Leading (algebraic) part}: $12 + 120 = 132$.
These are $1/|\zeta(-1)| + 1/|\zeta(-3)|$ --- values at
\emph{negative} integers, the non-perturbative wing. They
represent the ``vacuum contribution'' to $\alpha^{-1}$.

\item \textbf{Correction (transcendental) part}: $g \cdot 132$
and $\zeta(6) - \zeta(7)$. These involve periods at
\emph{positive} integers, the perturbative wing. They represent
the ``loop corrections'' to $\alpha^{-1}$.
\end{itemize}

\subsection{Deligne--Beilinson structure}

The Deligne period conjecture~\cite{Deligne1979} and Beilinson
conjectures~\cite{Beilinson1984} predict that critical $L$-values
of motives factor as
\begin{equation}
L^*(M, n) \sim c^+(M) \cdot \Omega(M)
\end{equation}
where $c^+$ is algebraic and $\Omega$ is a period.

Formula (\ref{eq:alpha}) has precisely this form:
\begin{equation}
\alpha^{-1} = \underbrace{132}_{\text{algebraic } c^+}
+ \underbrace{\text{transcendental corrections}}_{\text{period } \Omega}
\end{equation}

If this identification holds, $\alpha^{-1}$ is a physical instance
of the Beilinson conjectures, with the algebraic part from the
non-perturbative wing and the period part from the perturbative wing.

\section{Evidence: $\Omega_\Lambda$ combines both wings}

\subsection{The functional-equation dual form}

We have shown~\cite{wb-wdw} that the dark energy density parameter
\begin{equation}
\Omega_\Lambda = \frac{2\pi}{9}
\end{equation}
admits the representation
\begin{equation}\label{eq:omega-dual}
\Omega_\Lambda = \frac{16\,|\zeta(-1)|\,\zeta(2)}{\pi}
\end{equation}
which explicitly combines:
\begin{itemize}
\item $\zeta(-1) = -1/12$: non-perturbative wing (Bernoulli,
algebraic)
\item $\zeta(2) = \pi^2/6$: perturbative wing (period,
transcendental)
\item $1/\pi$: normalization (from Einstein's equation $8\pi G$)
\end{itemize}

\textbf{$\Omega_\Lambda$ lives on the bridge between the two
wings.}

\subsection{Wheeler--DeWitt derivation}

The derivation proceeds via a non-perturbative argument~\cite{wb-wdw}:

\begin{enumerate}
\item The Wheeler--DeWitt wave function $\Psi(a)$ satisfies temporal
Dirichlet--Neumann boundary conditions (Dirichlet at Big Bang,
Neumann at de Sitter asymptotic).
\item The 1D DN zero-point energy regularizes to
$\zeta_{\neg 2}(-1) = +1/12$ (non-perturbative wing, with Euler
$p = 2$ truncation).
\item The Cohen--Kaplan--Nelson holographic bound~\cite{CKN1999}
provides $\rho_\Lambda \leq M_P^2/\ell_H^2$ (the bridge between
UV and IR).
\item Saturation with coefficient $1/12$ yields
$\Omega_\Lambda = 2\pi/9$.
\end{enumerate}

Steps 1--2 are purely non-perturbative (no loops). Step 3 involves
the holographic principle (UV/IR mixing). The final result
(\ref{eq:omega-dual}) involves $\zeta(2)$ (perturbative wing) via
the $8\pi$ of Einstein's equation.

\subsection{Why perturbation theory cannot produce $\Omega_\Lambda$}

This explains why the cosmological constant problem has resisted
perturbative solution for 40+ years~\cite{Weinberg1989}. The
vacuum energy $\rho_\Lambda$ is a non-perturbative quantity
($\zeta$ at negative integers) that perturbation theory
($\zeta$ at positive integers) cannot access. Kreimer--Connes,
being perturbative, is structurally unable to predict
$\Omega_\Lambda$. The non-perturbative wing is required.

\section{The $p = 2$ Euler truncation as non-perturbative selection}

\subsection{DN boundary and $\zeta_{\neg 2}$}

The Euler product of $\zeta$ is
\begin{equation}
\zeta(s) = \prod_{p\text{ prime}} (1 - p^{-s})^{-1}.
\end{equation}
Removing the $p = 2$ factor gives
\begin{equation}
\zeta_{\neg 2}(s) = (1 - 2^{-s})\zeta(s)
= \prod_{p \geq 3} (1 - p^{-s})^{-1}.
\end{equation}

Physically, $\zeta_{\neg 2}$ arises from DN boundary conditions
(odd modes only). The removal of the $p = 2$ Euler factor
corresponds to the \emph{elimination of even harmonics} by the
mixed boundary.

\subsection{Invisibility to perturbation theory}

In perturbative QFT, all modes contribute to loop integrals.
There is no mechanism to select odd modes only. The Euler
product truncation at $p = 2$ is a \emph{boundary-condition
effect} on the vacuum, not a loop-level phenomenon.

\begin{proposition}[Non-perturbative character of $\zeta_{\neg 2}$]
The transition $\zeta \to \zeta_{\neg 2}$ (equivalently, DD $\to$
DN boundary) is invisible to perturbation theory. It is a
non-perturbative selection rule operating at the vacuum level.
\end{proposition}

This is why Kreimer--Connes never encounters $\zeta_{\neg 2}$.
The perturbative wing uses the \emph{full} $\zeta$ (all primes).
The non-perturbative wing uses $\zeta_{\neg 2}$ (odd primes only)
when DN boundaries are present.

\subsection{Sign consequence}

At $s = -1$:
\begin{align}
\zeta(-1) &= -\frac{1}{12} \quad \text{(NN/DD: attractive vacuum)} \\
\zeta_{\neg 2}(-1) &= +\frac{1}{12} \quad \text{(DN: repulsive vacuum)}
\end{align}

The sign flip is physically decisive: positive vacuum energy
($\rho_\Lambda > 0$) drives cosmic acceleration. Only DN boundary
conditions produce the correct sign. This is a non-perturbative
datum that perturbation theory has no access to.

\section{Unified picture}

\subsection{The zeta building}

We propose the metaphor of a ``zeta building'' with two wings:

\begin{equation}
\underbrace{\zeta(-5), \zeta(-3), \zeta(-1)}_{\text{Non-perturbative wing}}
\quad \xleftrightarrow[\text{functional equation}]{\zeta(s) \leftrightarrow \zeta(1-s)} \quad
\underbrace{\zeta(2), \zeta(3), \zeta(5), \zeta(7)}_{\text{Perturbative wing}}
\end{equation}

\begin{itemize}
\item \textbf{Left (non-perturbative) wing}: Bernoulli numbers,
vacuum zero-point, Casimir effect, $\Omega_\Lambda$,
boundary-condition dependence. Algebraic values. Accessible
via Wheeler--DeWitt quantum cosmology.

\item \textbf{Right (perturbative) wing}: Periods, MZVs, loop
integrals, anomalous dimensions, $g-2$. Transcendental values.
Accessible via Kreimer--Connes Hopf algebra.

\item \textbf{Connecting hallway}: The Riemann functional equation.
Physical observables like $\alpha^{-1}$ and $\Omega_\Lambda$ span
both wings.
\end{itemize}

\subsection{Physical observables and their wing structure}

\begin{center}
\begin{tabular}{llll}
\toprule
Observable & Left wing & Right wing & Bridge \\
\midrule
$a_e$ (electron $g-2$) & --- & $\zeta(3), \zeta(5)$ & KC only \\
Casimir (DD) & $\zeta(-3)$ & --- & Left only \\
Casimir (DN) & $\zeta_{\neg 2}(-3)$ & --- & Left only \\
$\Omega_\Lambda$ & $\zeta_{\neg 2}(-1)$ & $\zeta(2)$ via $8\pi G$ & Both \\
$\alpha^{-1}$ & $1/|\zeta(-1)|, 1/|\zeta(-3)|$ & $\zeta(6), \zeta(7)$ & Both \\
\bottomrule
\end{tabular}
\end{center}

Some observables live purely in one wing; others span both. The
spanning observables ($\alpha^{-1}$, $\Omega_\Lambda$) are the
most physically fundamental, as they require the full $\zeta$
building for their description.

\section{Relation to existing programs}

\subsection{Kreimer--Connes: the perturbative wing is established}

The right wing is mathematically rigorous and experimentally
verified (electron $g-2$ to $\sim 10^{-13}$ precision). Our
proposal does not modify or challenge Kreimer--Connes. It
\emph{extends} the same philosophy (``$\zeta$ values are
physically real'') to the non-perturbative side.

\subsection{Casimir physics: the left wing is partially established}

The Casimir effect for DD boundaries is experimentally
confirmed~\cite{Lamoreaux1997}. The $\zeta(-3) = 1/120$
coefficient appears in the measured force. Boyer's DN prediction
($\zeta_{\neg 2}(-3) = -7/120$) is not yet experimentally tested.

Our proposal extends this to cosmological scales via the
Wheeler--DeWitt temporal DN argument.

\subsection{Deligne--Beilinson: the bridge is conjectured}

The Deligne and Beilinson conjectures~\cite{Deligne1979,Beilinson1984}
predict that $L$-values factor into algebraic and transcendental
(period) parts. Our observation that $\alpha^{-1}$ and
$\Omega_\Lambda$ have precisely this structure is a candidate
physical realization.

\subsection{Moonshine: a parallel}

Monstrous Moonshine~\cite{Borcherds1992} demonstrated that
modular forms ($j$-function) encode genuine physics (Monster VOA,
2D CFT). Mathieu Moonshine~\cite{EOT2010} extended this to K3
surfaces. Our proposal that $\zeta$ special values encode 4D
physical constants is a parallel ``$\zeta$-moonshine'' --- not
via a sporadic group, but via the functional equation duality of
the $\zeta$ function itself.

\section{Predictions}

The unified framework makes specific testable predictions:

\begin{enumerate}
\item \textbf{$\Omega_\Lambda = 2\pi/9$}: Testable at $0.5\%$
by Euclid/DESI/LSST (2025--2028). This is a purely
non-perturbative-wing prediction, inaccessible to perturbation
theory.

\item \textbf{DN Casimir (Boyer)}: The 3D DN Casimir coefficient
$\zeta_{\neg 2}(-3) = -7/120$ (repulsive) can be tested in
laboratory experiments with mixed-boundary acoustic or
electromagnetic resonators.

\item \textbf{$\alpha^{-1}$ decomposition}: The split of
$\alpha^{-1}$ into algebraic (132) and transcendental
corrections can be compared against high-precision QED
calculations. If the split is exact, higher-loop terms should
organize into period-valued corrections to the Bernoulli
leading part.

\item \textbf{Muon and tau $g-2$}: The identification
$\Delta a_\mu \approx 252 = 6/|B_6|$ (non-perturbative wing)
and the Ramanujan discriminant $\tau(3) = 252$ suggest that
generational anomalies may bridge both wings. Prediction:
$\Delta a_\tau \propto \tau(5) = 4830$.
\end{enumerate}

\section{Limitations}

\begin{enumerate}
\item The ``two-wing'' picture is a \emph{proposal}, not a theorem.
The functional equation is a mathematical identity; its physical
interpretation as a duality between perturbative and
non-perturbative physics is conjectural.

\item The Wheeler--DeWitt derivation of $\Omega_\Lambda$ assumes
specific boundary conditions (temporal DN) and CKN bound
saturation. These assumptions are physically motivated but not
derived from first principles.

\item The decomposition of $\alpha^{-1}$ into algebraic + period
parts is observed post hoc. A rigorous derivation via
Kreimer--Connes Hopf algebra is future work.

\item The Euler $p = 2$ truncation (DN boundary) is presented as
a non-perturbative selection rule, but its fundamental origin
(why $p = 2$? why DN?) remains open. Candidates include spin
structure ($K_1(\mathbb{Z}) = \mathbb{Z}/2$) and time-reversal
asymmetry.

\item This framework does not address fermion masses, Yukawa
couplings, or CKM mixing --- these may require additional
structure beyond the $\zeta$ building.
\end{enumerate}

\section{Discussion}

The central thesis of this paper is that the arithmetic content
of quantum physics is organized by the Riemann $\zeta$ function,
with perturbative physics on one side of the functional equation
and non-perturbative vacuum physics on the other.

Kreimer--Connes brilliantly established the perturbative side.
They showed that $\zeta$ values at positive integers are not
regularization artifacts but genuine motivic periods with
Galois-theoretic structure. Our contribution is to extend this
philosophy to the other side: $\zeta$ values at negative integers
(Bernoulli numbers) appearing in vacuum energy regularization
are not computational conveniences but genuine physical data about
the boundary-condition structure of the vacuum.

If this picture is correct, the cosmological constant problem is
not a failure of quantum field theory per se, but a consequence
of \emph{looking only at the perturbative wing}. The full $\zeta$
building, spanning both wings via the functional equation,
naturally produces the observed small positive $\Lambda$ from
the non-perturbative wing ($\zeta_{\neg 2}(-1) = +1/12$) combined
with the holographic bridge (CKN bound).

The beauty of this picture lies in its economy: no new fields,
no extra dimensions, no multiverse. Only the recognition that
the $\zeta$ function's functional equation is not merely a
mathematical symmetry but a \emph{physical duality} between the
perturbative and non-perturbative sectors of quantum physics.

\section{Conclusion}

We have proposed that the physical content of the Riemann $\zeta$
function comprises two complementary wings:

\begin{itemize}
\item The \textbf{perturbative wing} ($\zeta$ at positive
integers: periods, MZVs, loop integrals), established by
Kreimer--Connes.
\item The \textbf{non-perturbative wing} ($\zeta$ at negative
integers: Bernoulli numbers, vacuum zero-point, boundary
conditions), proposed by this work via Wheeler--DeWitt quantum
cosmology.
\end{itemize}

The Riemann functional equation connects these wings, and specific
physical observables ($\alpha^{-1}$, $\Omega_\Lambda$) span both.
The framework provides:
\begin{enumerate}
\item A parameter-free prediction $\Omega_\Lambda = 2\pi/9$
(2\% agreement with Planck 2018).
\item A Deligne--Beilinson-type decomposition of $\alpha^{-1}$
into algebraic and period parts.
\item An explanation of why the cosmological constant problem
resists perturbative solution (it lives in the non-perturbative
wing).
\item A non-perturbative selection rule ($\zeta \to \zeta_{\neg 2}$,
DN boundary) invisible to loop expansion.
\end{enumerate}

The proposal positions vacuum cosmology as a natural complement
to Kreimer--Connes perturbative renormalization. Together, they
may constitute the two halves of a complete arithmetic description
of quantum physics --- united by the functional equation of the
Riemann $\zeta$ function.

\begin{thebibliography}{99}

\bibitem{Kreimer1998}
D.~Kreimer, Adv.\ Theor.\ Math.\ Phys.\ \textbf{2}, 303 (1998).

\bibitem{CK1998}
A.~Connes and D.~Kreimer,
Commun.\ Math.\ Phys.\ \textbf{199}, 203 (1998).

\bibitem{CK2000}
A.~Connes and D.~Kreimer,
Commun.\ Math.\ Phys.\ \textbf{210}, 249 (2000).

\bibitem{CK2001}
A.~Connes and D.~Kreimer,
Commun.\ Math.\ Phys.\ \textbf{216}, 215 (2001).

\bibitem{Brown2015}
F.~Brown, Ann.\ Math.\ \textbf{181}, 949 (2015).

\bibitem{BrownKreimer2013}
F.~Brown and D.~Kreimer,
Lett.\ Math.\ Phys.\ \textbf{103}, 933 (2013).

\bibitem{ConnesMarcolli2008}
A.~Connes and M.~Marcolli, \textit{Noncommutative Geometry,
Quantum Fields and Motives} (AMS, 2008).

\bibitem{Laporta2017}
S.~Laporta, Phys.\ Lett.\ B \textbf{772}, 232 (2017).

\bibitem{Aoyama2020}
T.~Aoyama \textit{et al.},
Phys.\ Rep.\ \textbf{887}, 1 (2020).

\bibitem{CKN1999}
A.~G.~Cohen, D.~B.~Kaplan, and A.~E.~Nelson,
Phys.\ Rev.\ Lett.\ \textbf{82}, 4971 (1999).

\bibitem{Casimir1948}
H.~B.~G.~Casimir,
Proc.\ Kon.\ Ned.\ Akad.\ Wet.\ \textbf{51}, 793 (1948).

\bibitem{Lamoreaux1997}
S.~K.~Lamoreaux,
Phys.\ Rev.\ Lett.\ \textbf{78}, 5 (1997).

\bibitem{Boyer1974}
T.~H.~Boyer, Phys.\ Rev.\ A \textbf{9}, 2078 (1974).

\bibitem{Weinberg1989}
S.~Weinberg, Rev.\ Mod.\ Phys.\ \textbf{61}, 1 (1989).

\bibitem{Deligne1979}
P.~Deligne, Proc.\ Symp.\ Pure Math.\ \textbf{33.2}, 313 (1979).

\bibitem{Beilinson1984}
A.~Beilinson, J.\ Soviet Math.\ \textbf{30}, 2036 (1985).

\bibitem{Borcherds1992}
R.~Borcherds, Invent.\ Math.\ \textbf{109}, 405 (1992).

\bibitem{EOT2010}
T.~Eguchi, H.~Ooguri, and Y.~Tachikawa,
Exp.\ Math.\ \textbf{20}, 91 (2011).

\bibitem{wb-alpha}
Wright Brothers Research Program,
``Physical constants as Eisenstein Fourier coefficients'' (2026).

\bibitem{wb-wdw}
Wright Brothers Research Program,
``Wheeler--DeWitt temporal DN vacuum and
$\Omega_\Lambda = 2\pi/9$'' (2026).

\end{thebibliography}

\end{document}
